\section{The fixed-index positive allowance at 39}
\label{ext:positive39-local}

The acquisition used by the sharp argument at $39$ requires the
intermediate cap $C_{39}<.41449$. It is sufficient first to exclude
\[
 C_{39}\ge\CE+K/39,\qquad K=741/4000,
\]
because $.40974+K/39=.41449$ and $\CE<.40974$.
This is a separate fixed-index comparison. Its inputs are
$C_{39}<.4226$, $39(C_{39}-\CE)<.50$, predecessor loss at most
$.33$, and both deleted score bounds at most $.4253$.
Sharp predecessors and the one-step estimate supply these inputs.
Indeed $\sqrt{39/38}<51/50$ and $\sqrt{38}>6$ give
\[
 \begin{gathered}
 39(C_{39}-\CE)<\frac{39}{38}
       \left(\frac{41}{200}+\frac{17}{250}+\frac{17}{1250}\right)
 =\frac{55887}{190000}<.3,\\
 C_{39}<.41+.3/39<.4226.
 \end{gathered}
\]
Under the contrary hypothesis the predecessor loss is zero, and the
two sharp deleted bounds are below $.4253$.
No positive-allowance result at another index is used.

The support polynomial in Section~\ref{support:polynomial} is
evaluated with
\[
             (N,L,K,q,R)=(39,39,741/4000,.33,3.84).
\]
Its derivation applies at this index, including the scale estimate
valid for $n\ge36$. Subsequent radius contractions use the same
fixed-index polynomial. Central bounded-threshold contractions use
the positive-error series with excess allowance $.50$ and loss $.33$.
The initializer covers all signed moment cells and all 100 probability
panels. It retains the lower score $C\ge.4095+K/39$ and the
charge $K\rho_l$, because $39(C-\CE)\rho\ge K\rho_l$.
This charge is part of the actual resource budget and is kept through
the coherent score corners and the later refinements.

The confinement proof of Section~\ref{conf:re:confinement} applies
with its actual parameters $N=39$, $J_0=37$ and deletion factor
strictly larger than $\sqrt{\max\{(39/38)^3,39/37\}}$.
The acquired initial strip is $(-w,1+w)$ with
$w\in\{.55,.575\}$; the conditional cubic coefficient is $30$
on the left and $24$ centrally. In the proof replace $.55$ by $w$
in the cubic bound, interval endpoints, width gates and upper
contradiction $L_g-wU_+>\mathcal E$. The gate
\[
 c_\rho^2[l(1-l)]^3>[3(1-l)+w]^2,
 \qquad [l,h]\subset(0,1/2),
\]
licenses the cubic bound. Together with $d_0<w$, the strip and the
inward endpoints give $|U_i|<w$. All positive slopes, inverse
denominators and signed budget gates are checked at the actual
index. After each contraction the newly proved amplitude is the
next strip. Thus this is the same confinement proof with explicit
parameters, without extending the fixed-index statement to another
integer.

The local comparison then consumes 78 probability rows covering
$[151/500,229/500]=[.302,.458]$, with all the probability and moment
complements excluded by the preceding acquisition. Its 624 signed
profile cells retain both signs of $H$ and four noise bands. The
whole-count part has 90 signed rows. Its closed phase cover is:
nearest-count comparison on $(0,.05]$ and $[.05,.08]$ at every
phase; nearest count for $|u|\ge.19$ and whole count for
$|u|\le.19$ on $[.08,.10]$; the corresponding boundary $.27$ on
$[.10,.12]$; and whole count at every phase above $.12$.
Each band is intersected with its own acquired noise ceiling.

Here the target is $C< C_B+K/39$, so the scalar tests differ from
the zero-allowance tests in Section~\ref{shortloc:section}.
All of them follow from the same identity
\eqref{shortloc:master}. Write $b=\sqrt{39}\pi_k$ and retain
the verified two-sided bound $m\le b\le1$. A nearest-count
bracket $D$ may now have either sign; consequently
\[
 bD\ge\min\{mD,D\}.
\]
If the repair is at most $r\lambda/39$, with $r\ge0$, a noise
band $[a,b_\lambda]$ is excluded by the intermediate-noise test
\[
 \min\{mD,D\}-r b_\lambda/39+K/39>0.
\]
Indeed the identity gives $\rho_B(C-C_B)<K/39$.
The elementary inequality $\rho_B\ge1$ follows from
$(1-2v)^2-v=(1-v)(1-4v)\ge0$, $0<v\le1/4$.
Thus $C<C_B+K/39\le\CE+K/39$.

At arbitrarily small positive noise the rescaled quartic and tail
estimates give the same inequality per unit $\lambda$. The tiny-noise
test adds the allowance
\[
             \frac{K}{39\max\{.05,\Lambda\}}
\]
to that per-noise reserve. Multiplying by
$0<\lambda\le\min\{.05,\Lambda\}$ again gives a bound
strictly below $K/39$. No positive lower noise cutoff is inserted.

For the whole-count test, write
\[
 S(\lambda)=\max\{mB(\lambda),B(\lambda)\}
                         +I(\lambda)/39.
\]
Both the bracket and the inverse correction may have either sign;
the first maximum uses $m\le b\le1$, while the division by $39$
uses the actual integer. Hence
$\rho_B(C-C_B)\le\lambda S(\lambda)$.
The saved gate on $[a,b_\lambda]$ is
\[
        K\rho_{B,l}/39-b_\lambda S(\lambda)>0,
\]
where $\rho_{B,l}$ is a positive lower bound for $\rho_B$.
If $S<0$ the conclusion is immediate; otherwise it follows by
$\lambda\le b_\lambda$ and $\rho_B\ge\rho_{B,l}$.
The formula is convex in positive $\lambda$, since its only
nonaffine terms are positive multiples of
$(\alpha\lambda+\beta)^{-1}$ and $1/\lambda$, and maxima of
convex functions preserve convexity. Both endpoint gates therefore
cover the whole noise cell.

Finally, zero noise makes the marginal exactly Bernoulli and gives
$C\le\CE<\CE+K/39$. The weak all-probability Bernoulli theorem
suffices; no finite binomial slack table is needed here. Once the
nine original arithmetic actions and their full source interfaces
are verified, these strict comparisons exclude the hypothesized
positive allowance and supply the cap required by the sharp39
initializer. The positive-allowance and sharp39 checks remain
distinct numerical obligations.
